\documentclass[10pt,a4paper]{amsart}
\usepackage[margin=19mm]{geometry}
\usepackage{amsmath,amssymb,mathtools}
\usepackage[scr=rsfs,cal=euler]{mathalfa}
\usepackage{xfrac}
\usepackage[unicode,hidelinks,bookmarks=false]{hyperref}
\newcommand{\ie}{i.e.\ }
\newcommand{\cf}{cf.\ }
\newcommand{\resp}{respectively\ }
\newcommand{\Subsection}{\subsection}
\providecommand{\nobreakdash}{\nobreak\hbox{-}}
\setlength{\emergencystretch}{2em}
\multlinegap0pt
\usepackage{bm}
\usepackage{bbm}
\usepackage{dsfont}
\usepackage{xspace}
\usepackage{cancel}
\usepackage{mathtools}
\usepackage{amsthm}
\makeatletter
\@ifundefined{defi}{\newtheorem{defi}{Defi}}{}
\@ifundefined{corollary}{\newtheorem{corollary}{Corollary}}{}
\@ifundefined{lemma}{\newtheorem{lemma}{Lemma}}{}
\@ifundefined{theorem}{\newtheorem{theorem}[defi]{\bf T\scriptsize HEOREM \normalsize}}{}
\makeatother
\makeatletter
\expandafter\ifx\csname corollary\endcsname\relax
\newenvironment{corollary}{\begin{dk}\it}{\end{dk}}
\fi
\expandafter\ifx\csname definition\endcsname\relax
\newenvironment{definition}{\begin{defi}\rm}{\end{defi}}
\fi
\newcommand{\deter}[1]{\det\zav{#1}}
\expandafter\ifx\csname example\endcsname\relax
\newenvironment{example}{\begin{exa}\rm}{$\kopr$\end{exa}}
\fi
\expandafter\ifx\csname lemma\endcsname\relax
\newenvironment{lemma}{\begin{lm}\it}{\end{lm}}
\fi
\expandafter\ifx\csname note\endcsname\relax
\newenvironment{note}{\begin{no}\rm}{\end{no}}
\fi
\expandafter\ifx\csname problem\endcsname\relax
\newenvironment{problem}{\begin{pro}\it}{\end{pro}}
\fi
\expandafter\ifx\csname proposition\endcsname\relax
\newenvironment{proposition}{\begin{prop}\it}{\end{prop}}
\fi
\expandafter\ifx\csname remark\endcsname\relax
\newenvironment{remark}{\begin{rem}\rm}{$\kopr$\end{rem}}
\fi
\newcommand{\tr}{{\rm tr}\,}
\newcommand{\trunc}[1]{\left\lfloor #1 \right\rfloor}
\newcommand{\zav}[1]{\left( #1 \right)}
\makeatother
\title[Coval description of the boundary of a numerical range and t]{Coval description of the boundary of a numerical range and the secondary values of a matrix}
\author{Petr Blaschke}
\date{Source: arXiv:2410.03744v3 (CC BY 4.0)}
\begin{document}
\maketitle

\noindent\emph{Source and licence.} Coval description of the boundary of a numerical range and the secondary values of a matrix. Version arXiv:2410.03744v3 (CC BY 4.0), distributed under the Creative Commons Attribution 4.0 International licence, as recorded in the arXiv licence metadata for this exact version and shown on the abstract page https://arxiv.org/abs/2410.03744v3. Full rights notice: licenses/arxiv-2410.03744-CC-BY-4.0.txt.\par\smallskip

\noindent\textit{Editorial scope.} Counted unit: the Corollary whose source label is \texttt{3by3cor} in the article (source0.tex lines 1579--1592, source label \texttt{3by3cor}), delivered together with the article's own complete original proof of it (source0.tex lines 1594--1665, one \texttt{proof} environment of 72 lines). No mathematics is written, repaired or re-derived: statement and proof are the article's own text line for line. Required context retained uncounted: maint (lines 1353--1363); anticharexplicit (lines 1488--1497). Every definition, display and label of those lines is retained unchanged. Omitted from this package and not redistributed: 1--1352, 1364--1487, 1498--1578, 1593--1593, 1666--2351. Nothing inside a retained excerpt is omitted or altered: every retained body is delivered exactly as the article's lines are written, so no omission receipt of any kind accompanies this package. Citations: the retained excerpts carry no citation command at all, so no bibliography entry of the article is transcribed and no reference is redistributed. Reference adaptation: none was needed and none was made; every reference inside the delivered unit resolves inside it, so no unresolved reference or citation is produced. Printed slips kept literal: no printed word, symbol, hypothesis or number of the counted statement or of its proof was altered. Layout and class: the article's own class is \texttt{amsart} with packages including inputenc, babel, color, hyperref, amssymb, amsmath, lmodern, fontenc, url, graphicx, standalone; the wrapper is the lane's pinned \texttt{amsart} editorial wrapper at 10pt on A4, which restates the theorem-like environments the retained text needs and adds the article's own macro definitions that the retained text needs (\texttt{\textbackslash deter}, \texttt{\textbackslash tr}, \texttt{\textbackslash trunc}, \texttt{\textbackslash zav}), with extra wrapper packages (bm, bbm, dsfont, xspace, cancel, mathtools). The delivered numbering is the unit's own. Assets: no retained span contains a figure environment, a graphics-inclusion command, a PDF-page inclusion, a table or a TikZ picture; no image file is copied into this package, the package declares no asset dependency and no image file is redistributed with it. Licence: version arXiv:2410.03744v3 of this article is distributed under the Creative Commons Attribution 4.0 International licence, as recorded in the arXiv licence metadata for this exact version and shown on the abstract page https://arxiv.org/abs/2410.03744v3. A full rights notice is delivered at licenses/arxiv-2410.03744-CC-BY-4.0.txt. Characters: every retained excerpt is pure ASCII, so the delivered engine, XeLaTeX with its default Latin Modern text font, reports no missing character.
\begin{theorem}\label{maint}
Let $A\in\mathbb{C}^{n\times n}$, $n\geq 2$. Then we have
\begin{equation}\label{covalpa}
\deter{p(A)}=\prod_{j=1}^{n}p(\lambda_j)-\frac{t_1(A)}{4}\prod_{j=1}^{n-2}p\zav{\lambda^{(2)}_j}+\sum_{k=2}^{\trunc{\frac{n}{2}}}\frac{\gamma_k}{4^k}\prod_{j=1}^{n-2k}p\zav{\lambda^{(k+1)}_{j}}.
\end{equation}
Here, $\gamma_k\in \mathbb{R}$,  $\lambda_j$ are the eigenvalues of $A$ and $\lambda^{(2)}_j$ are solutions of 
\begin{equation}\label{anticharacteristiceq}
\tr {\rm adj}\zav{x I-A}\zav{A^\star-\bar A}=0.
\end{equation}
The remaining poles might be at infinity, i.e. $\lambda_{j}^{(k)}\in \mathbb{C}^\star$ for $k>2$ in general. 
\end{theorem}

\begin{lemma}
Let $A\in\mathbb{C}^{n\times n}$ with eigenvalues $\lambda_1,\dots, \lambda_n$. It holds
\begin{equation}\label{anticharexplicit}
\tr {\rm adj}(xI-A)(A^\star -\bar A)=\sum_{j=1}^{n-1}\sum_{k=0}^{n-1-j}t_j(A)x^{n-1-j-k}(-1)^k e_k(A),
\end{equation}
where $e_k(A)$ are elementary symmetric polynomials
$$
e_k(A):=\sum_{1\leq j_1\leq j_2\leq\dots \leq j_k\leq n }\lambda_{j_1}\cdots \lambda_{j_n}.
$$
\end{lemma}

\begin{corollary}\label{3by3cor}
Let $A\in\mathbb{C}^{3\times 3}$ with $\lambda_1,\lambda_2,\lambda_3$ as its eigenvalues. Then the curve $\deter{p(A)}=0$ satisfies
\begin{equation}\label{covalpa3}
p(\lambda_1)p(\lambda_2)p(\lambda_3)=\frac{t_1(A)}{4}p(\chi),
\end{equation}
where the number $\chi$ is given by 
\begin{equation}\label{lambda2formula}
\chi:=\tr A-\frac{t_2(A)}{t_1(A)},\qquad t_1(A)\not=0.
\end{equation}
Also
\begin{equation}\label{chiinWA}
\chi\in W(A).
\end{equation}
\end{corollary}

\begin{proof}
The form of (\ref{covalpa3}) is the direct consequence of Theorem~\ref{maint}. The formula (\ref{lambda2formula}) follows from (\ref{anticharexplicit}).

For the second fact (\ref{chiinWA}), suppose we have a Schur form of $A$:
$$
A:=\zav{\begin{array}{ccc}
\lambda_1 & b_1 & b_2 \\
0 & \lambda_2 & c_1 \\
0 & 0 & \lambda_3 
\end{array}
},
$$
Then
$$
A^\star=\zav{\begin{array}{ccc}
\bar \lambda_1 & 0 & 0 \\
\bar b_1 & \bar \lambda_2 & 0 \\
\bar b_2 & \bar c_1 & \bar \lambda_3 
\end{array}
},
\qquad
\bar A=\zav{\begin{array}{ccc}
\bar \lambda_1 & * & * \\
0 & \bar \lambda_2 & * \\
0 & 0 & \bar \lambda_3 
\end{array}
},
$$
where stars $*$ represents some numbers that are irrelevant for the subsequent discussion.
Thus
$$
A^\star-\bar A =\zav{\begin{array}{ccc}
0 & * & * \\
\bar b_1 & 0 & * \\
\bar b_2 & \bar c_1 & 0 
\end{array}
},
$$
and
$$
t_1(A)=\tr A(A^\star-\bar A)=\tr \zav{\begin{array}{ccc}
\lambda_1 & b_1 & b_2 \\
0 & \lambda_2 & c_1 \\
0 & 0 & \lambda_3 
\end{array}
}\zav{\begin{array}{ccc}
0 & * & * \\
\bar b_1 & 0 & * \\
\bar b_2 & \bar c_1 & 0 
\end{array}
}=|b_1|^2+|b_2|^2+|c_1|^2.
$$
Also
\begin{align*}
t_2(A)&=\tr A^2(A^\star-\bar A)=\tr \zav{\begin{array}{ccc}
\lambda_1^2 & b_1(\lambda_1+\lambda_2) & b_2(\lambda_1+\lambda_3)+b_1c_1 \\
0 & \lambda_2^2 & c_1(\lambda_2+\lambda_3) \\
0 & 0 & \lambda_3^2 
\end{array}
}\zav{\begin{array}{ccc}
0 & * & * \\
\bar b_1 & 0 & * \\
\bar b_2 & \bar c_1 & 0 
\end{array}
}\\
&=|b_1|^2(\lambda_1+\lambda_2)+|b_2|^2(\lambda_1+\lambda_3)+|c_1|^2(\lambda_2+\lambda_3)+\bar b_2 b_1 c_1.
\end{align*}
Substituting these values for $t_1(A), t_2(A)$ into (\ref{lambda2formula}) yields
\begin{equation}\label{chiexplicit}
\chi=\frac{u A u^\star}{u u^\star},\qquad u:=(c_1,-b_2,b_1).
\end{equation}
\end{proof}

\end{document}
